\documentclass{plantillaPPT}

\titulo{10}{Coordenadas Polares}

\begin{document}
\primeraDiapo

\begin{frame}{Ecuaciones Polares}
\framesubtitle{Práctica 10}

1. Determinar la ecuación polar de las siguientes curvas:\\[4ex]

\((a)\) \(C:(x^2+y^2)^2 = 4(x^2-y^2)\) \\[4ex]
\((b)\) \(C:(x^2+y^2)^2 = x^2-y^2\) \\[4ex]
\((c)\) \(C: x^2-4y-4=0\)
    
\end{frame}

\begin{frame}{Integrales en coordenadas polares}
\framesubtitle{Práctica 10}

2. Sean \(R_1\) y \(R_2\) las regiones acotadas por las curvas \(C_1:x^2+y^2=8\), \(C_2:y=2\) y \(C_3: y=x\), y el eje \(x\), como se muestra en la figura adjunta:

\begin{center}
    \includegraphics[width=0.6\textwidth]{C2 2025-1/imagenes/ayud10img.png}
\end{center}
\end{frame}

\begin{frame}{Coordenadas polares}
\framesubtitle{Práctica 10}
\((a)\) Definir las regiones \(R_1\) y \(R_2\) de forma conjuntista en coordenadas rectangulares y polares.\\[5ex]
\((b)\) Escribir la o las integrales que permiten determinar el valor del área de la región \(R = R_1 \cup R_2\) en coordenadas polares.
    
\end{frame}

\begin{frame}{Integrales en coordenadas polares}
\framesubtitle{Práctica 10}

3. Sea \(r_1\) y \(r_2\) las curvas definidas por:
\[r_1=1 \qquad r_2=\sqrt{2\sin(2\theta)}, \quad \text{con }\theta\in[0,2\pi]\]

\begin{center}
    \includegraphics[width=0.5\textwidth]{C2 2025-1/imagenes/ayud10img2.png}
\end{center}

\end{frame}

\begin{frame}{Integrales en coordenadas polares}
\framesubtitle{Práctica 10}

\((a)\) Determinar los puntos de intersección entre ambas curvas.\\[3ex]
\((b)\) Calcular el área de la región interior a \(r_2\) y exterior a \(r_1\).\\[3ex]
\((c)\) Escribir la o las integrales que permitan calcular el área de la región interior a \(r_1\) y exterior a \(r_2\).
    
\end{frame}

\end{document}